% ======================================================================
% GENERAL DISCUSSION CHAPTER
% ----------------------------------------------------------------------
% I did not add a \begin{chapterabstractpage} here because this is the
% General Discussion chapter, unless you do want/need that, then by all means
% add it.% 
% ======================================================================
\section{General Discussion}

The Maastricht University PhD Thesis Template was developed to provide doctoral candidates with a reliable and elegant framework for writing, structuring, and publishing their dissertations in a professional and reproducible way. The template’s primary goal is to reduce the technical burden often associated with large LaTeX projects, allowing candidates to focus on the academic content rather than formatting concerns.

\vspace{1em}
From a design perspective, the template balances typographic quality and pragmatic usability. It integrates key features such as automated chapter title pages, customizable headers, and per-chapter abstracts, while maintaining full compatibility with Overleaf’s collaborative environment. By using LuaLaTeX or XeLaTeX, the template supports OpenType fonts and modern microtypography, ensuring visual consistency across operating systems and print outputs.

\vspace{1em}
Particular attention was given to transparency and modifiability. The accompanying template.sty file is fully commented, encouraging users to understand and adapt its components—geometry, typography, bibliographic style, and table of contents—according to their own faculty or university requirements. The structure is intentionally modular, with separate folders for chapters, front matter, and figures, reflecting good academic writing practice and version-control friendliness.

\vspace{1em}
While the current release defaults to APA referencing, the template allows seamless transitions to alternative citation styles such as Vancouver or Chicago. This flexibility ensures broad applicability across disciplines within Maastricht University and beyond. Furthermore, by including support for the university’s standard fonts (Thesis Sans and Verdana) but leaving the choice open to the user, the design respects both institutional identity and individual preference.

\vspace{1em}
PLEASE CHANGE THIS PLACEHOLDER TEXT WITH YOUR OWN GENERAL DISCUSSION.